\documentclass[10pt,a4paper]{article}

\usepackage[utf8]{inputenc}
\usepackage[T1]{fontenc}
\usepackage[french]{babel}
\usepackage{amsmath,amssymb,amsthm}
\usepackage{mathtools}
\usepackage{geometry}
\usepackage{graphicx}
\usepackage{hyperref}
\usepackage{booktabs}
\usepackage{algorithm}
\usepackage{algpseudocode}
\usepackage{xcolor}
\usepackage{tikz}
\usepackage{cite}

\geometry{margin=2cm}

\newtheorem{theorem}{Théorème}
\newtheorem{lemma}[theorem]{Lemme}
\newtheorem{proposition}[theorem]{Proposition}
\newtheorem{corollary}[theorem]{Corollaire}
\newtheorem{definition}{Définition}
\newtheorem{remark}{Remarque}

\title{\textbf{B-MoE : Mixture of Experts à Routage par Blocs et\\
Spécialisation Guidée des Experts}\\[0.5em]
\large Technical Report — Confidentiel}

\author{Sébastien Larnac}

\date{Juin 2026}

\begin{document}

\maketitle

\begin{abstract}
Nous proposons \textbf{B-MoE} (\textit{Bloc-routed Mixture of Experts}), une architecture de réseau de neurones profond qui étend les Mixture of Experts (MoE) standards selon trois axes complémentaires : (1) un routage inter-blocs opérant tous les $Z$ couches plutôt qu'à chaque couche, stabilisant le gradient et réduisant le coût de routage ; (2) un mécanisme de scoring basé sur une attention légère entre experts (\textit{inter-expert attention routing}), permettant une décision de routage informée par la sémantique collective des experts actifs ; (3) une initialisation guidée des experts par domaine (\textit{pre-specialization}), inspirée de la modularité fonctionnelle du cortex cérébral, accélérant la convergence et garantissant la diversité des représentations. Nous formalisons ces trois mécanismes, en établissons les propriétés théoriques principales, et discutons leur intégration dans un pipeline d'entraînement unifié.
\end{abstract}

\section{Introduction}

Les architectures Mixture of Experts \cite{shazeer2017outrageously} permettent d'augmenter massivement la capacité paramétrique d'un modèle sans augmentation proportionnelle du coût computationnel à l'inférence. Dans les formulations actuelles (Switch Transformer \cite{fedus2022switch}, Mixtral \cite{jiang2024mixtral}), le routage est effectué \textit{indépendamment à chaque couche}, sans mémoire de l'état des experts aux couches précédentes.

Cette indépendance couche-par-couche présente plusieurs limitations :
\begin{itemize}
  \item Le router à la couche $\ell$ ignore quels experts ont été activés aux couches $\ell-1, \ldots, \ell-k$, perdant ainsi une information de contexte précieuse.
  \item La multiplicité des décisions de routage introduit du bruit dans le gradient, rendant l'apprentissage des routers instable sur les réseaux profonds.
  \item Sans contrainte de spécialisation, les experts tendent vers la redondance (\textit{expert collapse}).
\end{itemize}

Nous proposons B-MoE pour répondre à ces trois limitations de manière unifiée.

\section{Formalisation du cadre}

\subsection{Notations}

Soit un transformer de profondeur $L$ composé de $N$ experts par position de routage. Notons :
\begin{itemize}
  \item $\mathbf{h}_\ell^{(t)} \in \mathbb{R}^d$ : le vecteur du \textit{residual stream} au token $t$ après la couche $\ell$
  \item $E_i^{(\ell)} : \mathbb{R}^d \to \mathbb{R}^d$ : la fonction de l'expert $i$ à la couche $\ell$
  \item $Z \in \mathbb{N}^*$ : la taille d'un bloc (nombre de couches entre deux décisions de routage)
  \item $B = \lfloor L/Z \rfloor$ : le nombre total de blocs
  \item $\sigma_b \in \{1, \ldots, N\}$ : l'indice de l'expert actif dans le bloc $b$
\end{itemize}

\subsection{Structure en blocs}

\begin{definition}[Bloc B-MoE]
Le bloc $b \in \{1, \ldots, B\}$ correspond aux couches $\ell \in \{(b-1)Z+1, \ldots, bZ\}$. Au sein d'un bloc, un unique expert $\sigma_b$ est activé pour l'ensemble des $Z$ couches. Le residual stream se propage selon :
\begin{equation}
  \mathbf{h}_\ell^{(t)} = \mathbf{h}_{\ell-1}^{(t)} + E_{\sigma_b}^{(\ell)}\!\left(\mathbf{h}_{\ell-1}^{(t)}\right), \quad \ell \in \text{bloc } b
\end{equation}
\end{definition}

La représentation agrégée à la sortie du bloc $b$ est définie par :
\begin{equation}
  \mathbf{o}_b^{(t)} = \mathbf{h}_{bZ}^{(t)}
\end{equation}

\section{Routage inter-blocs par attention}

\subsection{Limitation du routage scalaire standard}

Dans les MoE standards, le router est un MLP léger $r : \mathbb{R}^d \to \mathbb{R}^N$ qui produit un score scalaire par expert. Cette formulation ignore les relations \textit{entre} experts et ne permet pas au router de raisonner sur la complémentarité des spécialisations.

\subsection{Attention inter-experts}

\begin{definition}[Router par attention inter-experts]
Soient $\mathbf{o}_b^{(t)}$ les sorties des experts actifs au bloc $b$, indexées par $i \in \{1,\ldots,N\}$. On note $\mathbf{O}_b \in \mathbb{R}^{N \times d}$ la matrice des représentations agrégées de chaque expert (actif ou estimé). Le score de routage pour le bloc $b+1$ est calculé par :
\begin{equation}
  \mathbf{s}_{b+1} = \text{softmax}\!\left(\frac{\mathbf{Q}_b \mathbf{K}_b^\top}{\sqrt{d_k}}\right)\mathbf{V}_b \in \mathbb{R}^{N}
\end{equation}
où $\mathbf{Q}_b = \mathbf{O}_b W_Q$, $\mathbf{K}_b = \mathbf{O}_b W_K$, $\mathbf{V}_b = \mathbf{O}_b W_V$, avec $W_Q, W_K \in \mathbb{R}^{d \times d_k}$ et $W_V \in \mathbb{R}^{d \times 1}$ des matrices de projection apprises. L'expert sélectionné est :
\begin{equation}
  \sigma_{b+1} = \arg\max_{i}\, s_{b+1,i}
\end{equation}
\end{definition}

\begin{remark}
Le coût computationnel de ce router est $\mathcal{O}(N^2 d_k + N d)$ par token par décision de routage. Pour $N \leq 16$ et $d_k \ll d$, ce surcoût est négligeable devant le coût des couches transformer.
\end{remark}

\subsection{Propriété de différentiabilité}

\begin{proposition}[Différentiabilité du routage]
Le score $\mathbf{s}_{b+1}$ est différentiable par rapport aux paramètres $W_Q, W_K, W_V$ et aux représentations $\mathbf{O}_b$. La sélection $\arg\max$ est rendue différentiable via l'estimateur Straight-Through \cite{bengio2013estimating} ou Gumbel-Softmax \cite{jang2017categorical} lors de l'entraînement.
\end{proposition}

\begin{proof}
Par construction, $\mathbf{s}_{b+1}$ est une composition de fonctions différentiables (produit matriciel, softmax) en $W_Q, W_K, W_V$ et $\mathbf{O}_b$. Seule l'opération $\arg\max$ est discontinue. En substituant $\arg\max$ par $\text{softmax}(\mathbf{s}_{b+1}/\tau)$ avec température $\tau \to 0$ lors de l'entraînement (Gumbel-Softmax), ou en utilisant le gradient straight-through sur la sélection binaire, le graphe de calcul reste différentiable de bout en bout. $\square$
\end{proof}

\section{Pré-spécialisation des experts}

\subsection{Motivation neuroscientifique}

Dans le cortex cérébral humain, la modularité fonctionnelle est partiellement innée : l'aire de Broca présente une prédisposition structurelle au traitement du langage avant toute exposition linguistique \cite{dehaene2005neural}. Cette organisation n'est pas apprise \textit{ex nihilo} mais guidée par une initialisation génétiquement déterminée, accélérant considérablement la spécialisation post-natale.

Nous transposons ce principe aux MoE : au lieu d'initialiser tous les experts identiquement (ou aléatoirement), nous effectuons une \textit{pré-spécialisation} par domaine avant l'entraînement joint.

\subsection{Procédure de pré-spécialisation}

\begin{definition}[Corpus de pré-spécialisation]
Soit $\mathcal{D} = \{\mathcal{D}_1, \ldots, \mathcal{D}_N\}$ une partition de corpus par domaine (ex. : mathématiques, code, langage naturel, encyclopédique, multilingue, ...). Chaque expert $E_i$ est pré-entraîné sur $\mathcal{D}_i$ pendant $T_{\text{pre}}$ étapes via une loss de modélisation de langage standard :
\begin{equation}
  \mathcal{L}_{\text{pre}}(E_i) = -\mathbb{E}_{x \sim \mathcal{D}_i}\left[\log p_{E_i}(x)\right]
\end{equation}
\end{definition}

\begin{theorem}[Divergence garantie des experts]
Soient $E_i$ et $E_j$ deux experts pré-spécialisés sur des corpus $\mathcal{D}_i$ et $\mathcal{D}_j$ disjoints. Alors, sous des hypothèses de régularité standard (loss fortement convexe localement), les paramètres $\theta_i^*$ et $\theta_j^*$ issus de la minimisation de $\mathcal{L}_{\text{pre}}(E_i)$ et $\mathcal{L}_{\text{pre}}(E_j)$ vérifient :
\begin{equation}
  \|\theta_i^* - \theta_j^*\|_2 \geq \delta(\mathcal{D}_i, \mathcal{D}_j) > 0
\end{equation}
où $\delta(\mathcal{D}_i, \mathcal{D}_j)$ est une fonction croissante de la divergence KL entre les distributions $p_{\mathcal{D}_i}$ et $p_{\mathcal{D}_j}$.
\end{theorem}

\begin{proof}[Esquisse de preuve]
Par le théorème de convergence du gradient stochastique, $\theta_i^*$ est un minimiseur local de $\mathcal{L}_{\text{pre}}(E_i)$. Si $\mathcal{D}_i \cap \mathcal{D}_j = \emptyset$ et $D_{\mathrm{KL}}(p_{\mathcal{D}_i} \| p_{\mathcal{D}_j}) > 0$, alors les surfaces de loss $\mathcal{L}_{\text{pre}}(E_i)$ et $\mathcal{L}_{\text{pre}}(E_j)$ ont des minimiseurs distincts dans l'espace des paramètres. La borne inférieure sur $\|\theta_i^* - \theta_j^*\|_2$ découle de la $\mu$-forte convexité locale et de l'écart entre les gradients attendus sur $\mathcal{D}_i$ vs $\mathcal{D}_j$. $\square$
\end{proof}

\subsection{Entraînement joint avec régularisation de perméabilité}

Après pré-spécialisation, l'entraînement joint sur $\mathcal{D} = \bigcup_i \mathcal{D}_i$ doit maintenir la diversité des experts tout en permettant une spécialisation croisée partielle. On introduit une loss de régularisation :

\begin{equation}
  \mathcal{L}_{\text{div}} = -\lambda \sum_{i \neq j} D_{\mathrm{KL}}\!\left(p_{\theta_i} \| p_{\theta_j}\right)
\end{equation}

et une loss d'équilibrage de charge sur les trajectoires inter-blocs :

\begin{equation}
  \mathcal{L}_{\text{bal}} = \alpha \sum_{b=1}^{B} \sum_{i=1}^{N} \left(f_{b,i} - \frac{1}{N}\right)^2
\end{equation}

où $f_{b,i}$ est la fraction de tokens routés vers l'expert $i$ au bloc $b$.

La loss totale d'entraînement est :
\begin{equation}
  \mathcal{L} = \mathcal{L}_{\text{LM}} + \mathcal{L}_{\text{div}} + \mathcal{L}_{\text{bal}}
\end{equation}

\section{Propriétés du gradient}

\subsection{Stabilité sur les chaînes de routage}

\begin{proposition}[Réduction du vanishing gradient]
Dans un MoE à routage couche-par-couche, la norme du gradient de la loss par rapport aux paramètres du router à la couche $\ell$ décroît exponentiellement avec la profondeur. Dans B-MoE avec blocs de taille $Z$, le nombre de décisions de routage est réduit à $B = L/Z$, et le gradient du router au bloc $b$ vérifie :
\begin{equation}
  \left\|\frac{\partial \mathcal{L}}{\partial W_{\text{router}}^{(b)}}\right\| \geq \left\|\frac{\partial \mathcal{L}}{\partial W_{\text{router}}^{(\ell)}}\right\|_{\text{couche-par-couche}} \cdot C^{Z-1}
\end{equation}
pour une constante $C > 1$ liée à la norme spectrale des couches du bloc, toutes choses égales par ailleurs.
\end{proposition}

\begin{proof}
Dans le routage couche-par-couche, le gradient doit traverser $L$ opérations de routage discontinues. Avec le routage par blocs, il traverse $B = L/Z$ opérations. Le résiduel intra-bloc propage le gradient directement à travers $Z$ couches sans décision discrète intermédiaire, préservant la norme du gradient via les connexions résiduelles. $\square$
\end{proof}

\section{Complexité computationnelle}

\begin{proposition}[Complexité d'inférence]
Soit un transformer dense de référence de paramètres $P_{\text{dense}}$. Un B-MoE à $N$ experts, blocs de taille $Z$, avec $k=1$ expert actif par bloc vérifie :
\begin{align}
  P_{\text{total}} &= N \cdot P_{\text{dense}} \\
  P_{\text{actif}} &= P_{\text{dense}} + \mathcal{O}(N^2 d_k / Z) \approx P_{\text{dense}}
\end{align}
Le surcoût du router inter-experts est amorti sur $Z$ couches, soit un coût par couche de $\mathcal{O}(N^2 d_k / Z)$.
\end{proposition}

\section{Discussion}

\subsection{Comparaison avec les architectures existantes}

\begin{center}
\begin{tabular}{lccc}
\toprule
Architecture & Routage conditionnel & Spécialisation guidée & Stabilité gradient \\
\midrule
Switch Transformer & \texttimes & \texttimes & Faible \\
Mixtral (top-$k$) & \texttimes & \texttimes & Faible \\
Soft MoE & Partiel & \texttimes & Moyenne \\
Branch-Train-Merge & \texttimes & Oui (post-hoc) & N/A \\
\textbf{B-MoE (ce travail)} & \textbf{Oui} & \textbf{Oui (a priori)} & \textbf{Élevée} \\
\bottomrule
\end{tabular}
\end{center}

\subsection{Hyperparamètres critiques}

La valeur de $Z$ est le paramètre le plus sensible. Une valeur trop faible ($Z < 4$) réduit l'avantage sur le gradient sans apporter de bénéfice sémantique aux experts. Une valeur trop élevée ($Z > L/4$) limite la granularité du routage. Nous recommandons $Z \in [L/16, L/8]$ comme plage initiale de recherche, soit $Z \in [6, 12]$ pour $L = 96$.

\subsection{Limitations}

\begin{enumerate}
  \item La définition des corpus de pré-spécialisation $\mathcal{D}_i$ requiert une taxonomie préalable des domaines, qui peut être sous-optimale pour des tâches hybrides.
  \item L'hypothèse d'un seul expert actif par bloc ($k=1$) peut être insuffisante pour les tokens à sémantique composite. Une extension top-$k$ avec $k>1$ est directement applicable.
  \item L'estimation de $\mathbf{O}_b$ pour les experts non actifs au bloc $b$ requiert soit une évaluation forward partielle, soit une approximation par les représentations du bloc précédent.
\end{enumerate}

\section{Conclusion}

Nous avons proposé B-MoE, une architecture MoE qui intègre trois contributions orthogonales : routage par blocs pour la stabilité du gradient, attention inter-experts pour un scoring sémantiquement informé, et pré-spécialisation guidée pour la diversité et la convergence. Les propriétés théoriques établies — différentiabilité, divergence garantie des experts, réduction du vanishing gradient — constituent une base formelle solide pour une validation empirique sur des modèles de taille intermédiaire (1B–7B paramètres).

\bibliographystyle{plain}
\begin{thebibliography}{9}

\bibitem{shazeer2017outrageously}
N. Shazeer et al., \textit{Outrageously Large Neural Networks: The Sparsely-Gated Mixture-of-Experts Layer}, ICLR 2017.

\bibitem{fedus2022switch}
W. Fedus, B. Zoph, N. Shazeer, \textit{Switch Transformers: Scaling to Trillion Parameter Models with Simple and Efficient Sparsity}, JMLR 2022.

\bibitem{jiang2024mixtral}
A. Q. Jiang et al., \textit{Mixtral of Experts}, arXiv:2401.04088, 2024.

\bibitem{bengio2013estimating}
Y. Bengio, N. Léonard, A. Courville, \textit{Estimating or Propagating Gradients Through Stochastic Neurons for Conditional Computation}, arXiv:1308.3432, 2013.

\bibitem{jang2017categorical}
E. Jang, S. Gu, B. Poole, \textit{Categorical Reparameterization with Gumbel-Softmax}, ICLR 2017.

\bibitem{dehaene2005neural}
S. Dehaene, L. Cohen, \textit{The Unique Role of the Visual Word Form Area in Reading}, Trends in Cognitive Sciences, 2005.

\end{thebibliography}

\end{document}
